\subsection{Further attempts at the two remaining torus questions}

\statement{C.24}{Proved}{The proposed cubic lattice admits an invariant unimodular surface form}
Let $q\equiv1\pmod3$ be an odd prime.  The
$\mathbb Z_3[C_3]$-module
\[
 L_q=\mathbb Z_3[\zeta_3]\oplus
                       \mathbb Z_3[C_3]^{\,q-1}             \tag{C.24.1}
\]
admits a $C_3$-invariant unimodular alternating
form.  All integral fixed vectors have pairings
divisible by three.  This module is realized
by the homology of a geometric order-three
action on a closed surface of genus
$(3q-1)/2$.  Thus a contradiction based solely
on the module (C.24.1), its surface form and
mod-three isotropy of integral fixed characters
cannot exclude a scalar cubic torus coordinate.
This is an algebraic limitation of that
approach, not a realization of such a torus
coordinate.

\begin{proof}
On the cyclotomic summand take
\[
 T=\begin{pmatrix}0&-1\\1&-1\end{pmatrix},
 \qquad J=\begin{pmatrix}0&1\\-1&0\end{pmatrix}.
\]
Then $T^3=I$, $T^2+T+I=0$, and
$T^{\mathsf T}JT=J$.  This summand has no
nonzero integral fixed vector.
Since $q-1$ is even, pair the regular modules
in $(q-1)/2$ pairs.  On a pair with bases
$e_0,e_1,e_2$ and $f_0,f_1,f_2$, let the
generator cyclically permute each basis and
put $\langle e_i,f_j\rangle=\delta_{ij}$,
with the reverse pairing negative and all
other pairings zero.  This form is invariant
and unimodular.  The fixed submodule is
generated by $e_0+e_1+e_2$ and
$f_0+f_1+f_2$; their pairing is three.
Orthogonal direct sum proves both lattice
assertions.

For the surface realization, begin with an
elliptic curve admitting multiplication by
$\zeta_3$, whose order-three action has three
fixed points and cyclotomic integral homology.
Choose a free orbit of three small discs and
attach three handles equivariantly.  This
adds one pair of regular homology modules:
the three handle $a$ cycles and the three
handle $b$ cycles are cyclically permuted and
have the displayed hyperbolic pairings.
Repeat $(q-1)/2$ times, using disjoint free
orbits.  Genus becomes
$1+3(q-1)/2=(3q-1)/2$, and homology is
(C.24.1).  The three original fixed points
remain.  This construction has global
orientation multiplier one and does not
descend to the hypothetical scalar torus
action.  It proves exactly the asserted
compatibility of the linear data.
\end{proof}

\statement{C.25}{Proved}{A free-factor route would settle the actual fixed kernel}
In the notation of C.12, if the actual other
fixed group $B$ is a free profinite factor
of $F=\widehat F_2$, then $B\cong\widehat{\mathbb Z}$
and $K=\ker(\chi|_B)=1$.  The abstract
projective retract constructed in C.2 is
not a free profinite factor of $\widehat F_2$.

\begin{proof}
Suppose $F=B\amalg D$.  Both factors are
retracts of the finitely generated group
$F$ and are therefore finitely generated.
On abelianization this is the direct sum
$F^{\rm ab}=B^{\rm ab}\oplus D^{\rm ab}$.
The inclusion of $B^{\rm ab}$ is injective
and lies in the fixed module of the
negative involution, since $B$ is its
actual fixed subgroup.  At each prime this
fixed module has rank one: the order-two
matrix has determinant $-1$ and eigenvalues
$1,-1$.  The epimorphism
$\chi:B\to\widehat{\mathbb Z}$ forces
$B^{\rm ab}$ to have rank at least one
at every prime.  As a direct summand of
$\widehat{\mathbb Z}^2$, it is consequently
$\widehat{\mathbb Z}$, and so is
$D^{\rm ab}$.

For every finite group $A$, the universal
property of the full profinite free product
gives the exact count
\[
 |\operatorname{Hom}(B,A)|
 |\operatorname{Hom}(D,A)|=|A|^2.                         \tag{C.25.1}
\]
Each factor on the left is at least $|A|$,
because homomorphisms through its abelianization
$\widehat{\mathbb Z}$ are parametrized by the
image of one topological generator.
A homomorphism of $B$ with nonabelian image
would add another homomorphism beyond these
$|A|$, contradicting (C.25.1).  Hence every
finite image of $B$ is abelian, so $B$ is
abelian and C.1 gives $K=1$.

For the C.2 model the abelianization is also
$\widehat{\mathbb Z}$, but it has twelve
homomorphisms to $S_3$.  To see this, its
pro-three generators $x,y$ map into $A_3$.
The prime-to-three generator $z$ maps to
identity or a transposition.  If it is
identity, $zyz^{-1}=y^{-1}$ forces $y=1$
and gives three choices for $x$.
For each of its three transposition values,
$zxz^{-1}=x$ forces $x=1$ and gives three
choices for $y$.  Thus there are
$3+3\cdot3=12$ maps.  A free complement
would have procyclic abelianization and
give at least six maps to $S_3$, contradicting
$12\cdot6>6^2$.  The finite enumeration in
\path{code/check_cubic_and_factor.py} verifies
this count and the lattice matrices in C.24.
\end{proof}

\statement{C.26}{Open}{The scalar cubic existence question in prompt C.3}
We have not constructed or excluded an
order-three element of the full torus
loop-outer group whose multiplier has a
nonidentity cubic coordinate.  C.16 is
specific to the three-cusp structure and
does not prove the torus assertion.
C.24 explains exactly why replacing the
quartic involution in C.22 by the proposed
cubic integral module does not yield the
same contradiction: the fixed vectors
come from norms in regular modules and
their pairings vanish modulo three.
The module is even geometrically realizable
on the filled covering surface, while the
needed descent to the original punctured
torus with its specified nontrivial multiplier
is absent.

An additional cyclic-cover attempt has the
following precise limit.  At a scalar cubic
coordinate $q$, an eigencharacter
$F\to C_q$ is preserved up to a primitive
cube-root scalar, rather than pointwise.
The induced order-three action consequently
permutes nontrivial deck eigenspaces in
orbits of length three.  Its commutant is
no longer the product of unramified fields
with multiplicity one used in C.16; the
resulting permutation blocks themselves
have order three.  Thus absence of a
nontrivial cubic root in an unramified
triadic field does not exclude this different
semilinear action.
The finite Frobenius quotients proposed in
LM4 D.71 and D.73 would impose stronger
full-group constraints; without a separate
proof that those constraints are incompatible,
or a construction of a compatible full
action, they do not answer this row.

\statement{C.27}{Proved}{Mixed-sign diagonal power involutions fail the full torus cusp test}
For $u\in\widehat{\mathbb Z}^{\times}$ with $u^2=1$, the
continuous automorphism
\[
 \alpha_u:\widehat F(x,y)\longrightarrow\widehat F(x,y),
 \qquad x\longmapsto x^u,\quad y\longmapsto y^{-u}           \tag{C.27.1}
\]
has order two and determinant $-1$ on full abelianization.
It preserves the torus cusp $[x,y]$ up to conjugacy and
the power $-1$ if and only if $u=1$ or $u=-1$ globally.
In these two cases it is a standard reflection and its
noncusp fixed group is $\widehat{\mathbb Z}$.

Every pro-$p$ image of (C.27.1) is a standard reflection,
including when the full cusp assertion fails.  Thus this
family gives an explicit mixed-prime obstruction to
one natural attempted construction in prompt C.2.
It does not classify general negative torus involutions.

\begin{proof}
Each primary coordinate $u_p$ equals $1$ or $-1$, since
$\mathbb Z_p$ is an integral domain.  The universal
property of the full free group defines (C.27.1), and
applying it twice fixes both generators.  Its determinant
is $-u^2=-1$, so it is not the identity.  At any one prime
it fixes one generator and inverts the other.  This
proves the assertions about the actual automorphism and
its pro-$p$ images.

Suppose first that the signs at two odd primes differ.
Name them $q,p$ so that $u_q=1$ and $u_p=-1$.  Put
$d=\operatorname{ord}_p(q)$, and choose an element
$\zeta$ of order $p$ in $\mathbb F_{q^d}^{\times}$.
Its degree over $\mathbb F_q$ is $d$.  Form the finite
group
\[
 G=(\mathbb F_{q^d},+)\rtimes C_p,
 \qquad (v,i)(w,j)=(v+\zeta^i w,i+j).
                                                               \tag{C.27.2}
\]
The elements $X=(1,0)$ and $Y=(0,1)$ generate $G$:
the conjugates of $X$ span the field additively because
$\zeta$ has degree $d$.  They have orders $q,p$,
respectively.  The epimorphism $x\mapsto X,y\mapsto Y$
therefore sends both $\alpha_u(x)$ and $\alpha_u(y)$
to the original generator images.  However
\[
 [X,Y]=(1-\zeta,0)
\]
is not conjugate to its inverse.  Conjugation on this
nonzero additive element multiplies it by a power of
$\zeta$, whereas $-1$ is not a power of an element of
odd order.  Thus the required full cusp conjugacy is
impossible in this finite quotient.

It remains to exclude a common sign at all odd primes
with the opposite sign at two.  In $\operatorname{SL}_2
(\mathbb F_7)$ take
\[
 X=\begin{pmatrix}0&6\\1&3\end{pmatrix},\qquad
 Y=\begin{pmatrix}0&2\\3&6\end{pmatrix},\qquad
 C=[X,Y]=\begin{pmatrix}5&6\\1&0\end{pmatrix}.              \tag{C.27.3}
\]
They have orders $8,3,14$, respectively.  In particular
$X^4=C^7=-I$.  The element $-C=I+N$ has $N\ne0$,
$N^2=0$.  Conjugacy of $C$ to $C^{-1}$ would conjugate
$N$ to $-N$.  In a basis adapted to the kernel of
$N$, a determinant-one conjugator has diagonal entries
$a,a^{-1}$, and changes the nonzero upper-right entry
of $N$ by $a^2$.  It would follow that $a^2=-1$ in
$\mathbb F_7$, which is false.  Hence $C$ is not
conjugate to $C^{-1}$ in this finite group.

If $u_2=1$ and all odd signs are $-1$, (C.27.1) fixes
both matrices in (C.27.3).  If $u_2=-1$ and all odd
signs are $1$, it inverts both.  In the latter case
\[
 [X^{-1},Y^{-1}]=(XY)^{-1}[X,Y](XY),
\]
so the commutator image still lies in the conjugacy
class of $C$, not that of its inverse.  Both remaining
mixed patterns are therefore excluded.

For $u=1$, the identity
$[x,y^{-1}]=y^{-1}[x,y]^{-1}y$ proves cusp preservation;
for $u=-1$ use
$[x^{-1},y]=x^{-1}[x,y]^{-1}x$.
These reflections are conjugate by the basis exchange.
For the first one, with section generator $s$, setting
$t=ys$ gives the full presentation
\[
 \widehat F_2\rtimes_{\alpha_1} C_2
       =(\widehat{\mathbb Z}\times C_2)\amalg C_2,
 \qquad [x,s]=1,\quad s^2=t^2=1.
\]
The centralizer theorem for profinite free products,
\cite[Theorem B']{A_HR1985}, identifies its noncusp
geometric centralizer with $\overline{\langle x\rangle}$.
The other case exchanges $x,y$.  Thus $K'=1$ for
this family whenever it is actually peripheral.

The executable \path{code/check_diagonal_involutions.py}
checks the order-$21$ instance of (C.27.2), enumerates
all $336$ matrices of the group in (C.27.3), verifies
that $X,Y$ generate it, and checks the two disjoint
commutator conjugacy classes.  The general argument
above supplies the assertion for every sign pattern.
\end{proof}
